%%%PAGE 258 rot=0 legibility=good summary=圆周运动基本量与向心力公式；传动装置（皮带/齿轮/链条/同轴/摩擦传动）；力学两确定；动态平衡解法提纲
%%%BLOCK id=258-01 ch=04 sec=圆周运动·基本量 kind=formula alt=none cont=none y=0.07-0.22
\begin{gather*}
v=\frac{\Delta\widehat{l}}{\Delta t}\qquad \unit{m/s}\\
\omega=\frac{\Delta\theta}{\Delta t}\ \ \text{矢量}\qquad \unit{rad/s}\\
n=f=\frac{\omega}{2\pi}\qquad \unit{r/s}
\end{gather*}
匀圆\quad $a_{\text{合}}=a_n$\quad $F_{\text{合}}=F_n$
\[ F_n=\frac{mv^2}{r}=m\omega^2 r=mv\omega=ma=m\frac{4\pi^2}{T^2}r \]
%%%END
%%%BLOCK id=258-02 ch=04 sec=皮带齿轮·传动装置 kind=knowledge alt=none cont=none y=0.17-0.40
\textbf{传动装置}：线速度相等
\begin{itemize}
\item 皮带（静摩擦力）
\item 齿轮传动\quad $T_{\text{分}}=60T_{\text{秒}}$\quad \unclear{$T_{\text{分}}=60\,r_{\text{秒}}$针$+$}
\item $f_{B\text{带}}$\ 链条
\end{itemize}
自行车\ 主\struck{\unclear{…}}动轮\quad 受地\unclear{$f$}向前
%FIG p258_a 0.565 0.225 0.762 0.295
\notefig[0.3]{figures/p258_a.png}
\[ r_A=\frac{1}{2}r_B,\qquad V_A=V_B,\qquad \frac{\omega_A}{\omega_B}=2 \]
受主动轮\unclear{转}方向\par
与主动轮转向一样
%FIG p258_b 0.585 0.300 0.755 0.350
\notefig[0.3]{figures/p258_b.png}
%%%END
%%%BLOCK id=258-03 ch=04 sec=皮带齿轮·传动装置 kind=knowledge alt=none cont=none y=0.37-0.52
\textbf{角速度相等}\ —\ 同轴传动
%FIG p258_c 0.12 0.405 0.285 0.472
\notefig[0.25]{figures/p258_c.png}
\[ r_B=3r_A,\qquad V_B=3V_A \]
%%%END
%%%BLOCK id=258-04 ch=04 sec=皮带齿轮·传动装置 kind=knowledge alt=none cont=none y=0.38-0.54
%FIG p258_d 0.355 0.390 0.762 0.535
\notefig[0.7]{figures/p258_d.png}
\[ V_1=V_2,\qquad \omega r\ \text{相等},\qquad n_1D_1=n_2D_2 \]
可动轴；主；越转越快
\[ r_1\omega_1=r_2\omega_2 \]
主；从；可变半径
%%%END
%%%BLOCK id=258-05 ch=02 sec=受力分析·力学两确定 kind=method alt=03 cont=none y=0.55-0.62
\textbf{力学两确定}
\begin{itemize}
\item 状态
  \begin{itemize}
  \item 平衡
  \item 非平衡
  \end{itemize}
\item 研究对象
\end{itemize}
%%%END
%%%BLOCK id=258-06 ch=02 sec=共点力平衡·动态平衡 kind=summary alt=none cont=none y=0.62-0.88
\textbf{动态平衡}
% 原稿：前三项另有一小括号相连，七项再由一个大括号统领
\begin{itemize}
\item 力三角形、平移、延长\ 旋转
\item 正弦、余弦定理、\unclear{余力}勾股定理\ 三角函数
\item 非特殊三角形\ 用相似
\item 定角画圆
\item 晾衣\unclear{架}模型\ $+$\ 转动 % 「架」处原稿字迹有叠写
\item 正交分解
\item 摩擦角、全反力
\end{itemize}
%%%END
%%%PAGEEND 258
